\chapter*{An Arithmetic Field Guide}
\addcontentsline{toc}{chapter}{An Arithmetic Field Guide}
\markboth{An Arithmetic Field Guide}{An Arithmetic Field Guide}

The first move on an arithmetic problem is to make its question smaller.
If it asks for a remainder, do not compute an enormous integer before
reducing the factors. If it asks for an integer solution, check small
moduli before guessing a formula. If it asks for all prime factors of
a number, distinguish that from merely deciding whether it is prime.

Here is a short route through three of the book's central ideas.

\begin{examplex}[from a gcd to a modular inverse]
The equation $35x\equiv1\pmod{12}$ asks for an integer $x$ such that
$35x-1$ is a multiple of $12$. First divide: $35=2\cdot12+11$ and
$12=11+1$. Reading backwards gives
\[
1=12-11=3\cdot12-35.
\]
Modulo $12$ this says $(-1)\cdot35\equiv1$, so
$x\equiv-1\equiv11\pmod{12}$. Check:
$35\cdot11=385=32\cdot12+1$.

Now replace $35$ by $18$. Since $\gcd(18,12)=6$,
every integer combination $18x+12y$ is a multiple of $6$.
It cannot equal $1$. Therefore $18x\equiv1\pmod{12}$
has no solution. The difference between the two equations is
coprimality, not the size of the coefficient.
\end{examplex}

The same test repeats throughout the volume. Before cancelling a factor modulo $m$,
check its gcd with $m$. Before factoring, name the ring. Before treating
a density as a probability, state the finite experiment and the
limit that connects them. The check at the end is part of the
argument, not a substitute for it.

\section*{Notation at a glance}
We use $\N=\{0,1,2,\ldots\}$ and $\N^*=\{1,2,3,\ldots\}$. A prime is a positive integer greater than $1$
with exactly two positive divisors. Thus $1$ is a unit, not a prime; $2$ is a prime, not an inconvenient exception
to be forgotten. Unless another domain is named, divisibility and greatest common divisors concern integers,
and $\gcd(a,b)$ is positive when $(a,b)\ne(0,0)$; $\gcd(0,0)=0$.

\begin{description}
\item[$a\mid b$.] There is an integer $c$ with $b=ac$. This is a statement, not a fraction.
\item[$a\equiv b\pmod m$.] $m\mid a-b$, with $m\ge1$. A residue class is a set of integers; a representative is
one integer chosen from it.
\item[$(\Z/m\Z)^\times$.] The invertible residue classes. They are precisely the classes coprime to $m$;
there are $\varphi(m)$ of them (Proposition~\ref{prop:units}).
\item[$\operatorname{ord}_m(a)$.] For $\gcd(a,m)=1$, the least positive exponent $r$ for which
$a^r\equiv1\pmod m$. It divides $\varphi(m)$; equality means that $a$ is a primitive root.
\item[$v_p(n)$.] The exponent of the prime $p$ in a non-zero integer $n$. For non-zero rationals,
$v_p(a/b)=v_p(a)-v_p(b)$. The valuation of zero and the corresponding $p$-adic absolute value belong to the optional later-volume preview below.
\item[$\tau(n),\ \sigma(n)$.] For an integer $n\ge1$, the number and the sum of its positive divisors, respectively.
\item[$\mu(n)$.] $0$ if a prime square divides $n$, otherwise $(-1)^r$ if $n$ is a product of $r$ distinct
primes; in particular $\mu(1)=1$.
\item[$\left(\frac ap\right)$.] The Legendre symbol for an odd prime $p$: $0$ if $p\mid a$, $1$ if $a$ is a
non-zero square modulo $p$, and $-1$ otherwise. A Jacobi symbol with composite denominator is not the same test.
\item[{$[a_0;a_1,a_2,\ldots]$.}] A simple continued fraction. The convergent $p_k/q_k$ uses the terms through
$a_k$. The subscript starts at $0$.
\item[$N\alpha$.] The norm of a Gaussian integer $\alpha=a+bi$ is
$a^2+b^2$. The Eisenstein norm likewise comes from multiplying an element
by its complex conjugate; Chapter~\ref{ch:gaussian} proves the corresponding
division and factorisation results.
\item[$\chi$, $L(s,\chi)$.] A Dirichlet character modulo $m$ (Definition~\ref{def:t1-character}), and for $s>1$ its Dirichlet series $\sum\chi(n)n^{-s}$. The principal character is $\chi_0$. The function $\chi_4$ is the non-principal character modulo $4$.
\item[$\Phi(x,z)$.] For real $x\ge0$ and $z\ge2$, the number of integers $1\le n\le x$ with no prime factor $\le z$. Theorem~\ref{thm:t1-sieve} writes it as a M\"obius sum of floors.
\end{description}

\section*{Which tool should I try?}
\begin{description}
\item[A divisibility statement.] Apply Euclidean division; find a B\'ezout identity; or compare $v_p$ on each
side for every prime. Check whether the argument needs coprimality.
\item[A congruence.] Reduce first. Before cancelling $c$ modulo $m$, compute $\gcd(c,m)$. Separate coprime prime
powers with the Chinese remainder theorem (Theorem~\ref{thm:crt}).
\item[A very large exponent.] Check the base is a unit, then reduce the exponent modulo its order. Repeated
squaring works even when Euler's theorem does not.
\item[A sum over divisors.] Write the sum over pairs $ab=n$. Convolution or M\"obius inversion
(Theorem~\ref{thm:mobius}) may undo a complicated-looking operation.
\item[A square modulo a prime.] Factor the numerator, extract squares, and use the supplementary laws and
quadratic reciprocity (Theorem~\ref{thm:qr}). Keep the reciprocity sign visible.
\item[An integer equation.] Look modulo a small prime or prime power before searching. Then try a factorisation,
a valuation, a parametrisation, or infinite descent. A search proves completeness only after a bound reduces all possible candidates
to a finite list.
\item[A good approximation.] Use convergents and their determinant identity. For Pell's equation, retain exact
integer arithmetic: a small floating-point error is not a small error in an equation demanding equality to $1$.
\item[Broken unique factorisation.] Identify units, norms and irreducibles in the actual ring. An irreducible need
not be prime. In $\Z[\sqrt{-5}]$ the two factorisations of $6$ are the example in the Gaussian chapter; Volume~13 passes to ideals and restores unique factorisation.
\item[A root modulo $p^k$.] Check the derivative before lifting. A simple root lifts uniquely; a singular root
may branch or die. This lifting argument is an optional extension of the congruences chapter.
\item[A residue class of primes.] Test a character, or a cyclotomic polynomial, before reaching for Dirichlet's theorem. The classes modulo $4$ are separated by $\chi_4$. The exact sifted count is Theorem~\ref{thm:t1-sieve}; an asymptotic count is a later theorem.
\item[Squares modulo a power of $2$.] For $k\ge3$, reduce the unit modulo $8$. It is a square modulo $2^k$ if and only if it is $1$ modulo $8$ (Corollary~\ref{cor:mod8squares}).
\end{description}

\section*{Five questions before calling a proof finished}
\begin{enumerate}
\item What are the variables' domains? In particular, is the modulus prime, is the denominator non-zero,
and are the quantities positive where an order argument is used?
\item Have I proved existence, uniqueness, or both? Finding one root does not list all roots.
\item Have I used an implication backwards? If every square passes a test, passing that test need not make a
number a square. The Jacobi symbol supplies memorable counterexamples.
\item Does an iteration stop, or does it converge? An integer descent needs a strictly decreasing positive
integer; a limiting argument needs a specified complete space.
\item Did I prove the claim for every allowed input, including boundary cases, or merely check many examples?
Experiments are wonderful scouts. They cannot sign the treaty.
\end{enumerate}

\section*{One investigation, from evidence to a theorem}
Start with $f(n)=n^2+1$. Its values at $1,2$ are primes, $2,5$.
The next value, $10$, defeats the conjecture \emph{always prime}.
A more useful experiment records prime factors:
$10=2\cdot5$, $26=2\cdot13$, $50=2\cdot5^2$.
The odd factors seen here are all $1$ modulo $4$.
That suggests a different statement, about factors rather than values.

State it precisely: for every integer $n$ and odd prime $p$,
if $p\mid n^2+1$, then $p\equiv1\pmod4$.
Reduction asks whether $-1$ is a square; Euler's criterion proves the claim,
as in Exercise~\ref{ex:onemod4}. The exceptional prime $2$ is visible in the
data and must remain excluded. A failed conjecture has uncovered a theorem.

Change the question again. For each prime $p\equiv1\pmod4$,
Euler's criterion supplies a root $a$ of $X^2+1$ modulo $p$;
then $p\mid(a+kp)^2+1$ for every integer $k$.
There are infinitely many such primes by the same exercise.
Neither fact proves that $n^2+1$ itself is prime infinitely often,
an open problem. The quantifiers and the word \emph{factor} changed;
the proof must not silently change with them.

A useful investigation record is short: examples, the first failed
conjecture, a precise replacement, the proof's essential reduction,
and its exceptions. The Gaussian factorization
$n^2+1=(n+i)(n-i)$ offers another language for the same pattern.
Changing languages should explain the original question.

\section*{Optional later-volume preview: ideals and $p$-adic notation}
The following notation belongs to later work on general number fields and
$p$-adic numbers. It is not required to solve problems on the first route
through Volume~2. The Gaussian and Eisenstein rings in Chapter~\ref{ch:gaussian}
are concrete examples, not a substitute for the general theory.

\begin{description}
\item[$v_p(0)$.] In the extended valuation convention, $v_p(0)=+\infty$. This is useful when building the $p$-adic absolute value; it is not part of the first-route divisibility rules for non-zero integers.
\item[$\mathcal O_K$.] The ring of algebraic integers in a number field $K$. An algebraic integer
is a root of a monic polynomial with integer coefficients; a number field
is a field containing $\Q$ of finite dimension as a $\Q$-vector space.
For $K=\Q(i)$ and $K=\Q(\omega)$, the full rings are respectively
$\Z[i]$ and $\Z[\omega]$. Volume~13 proves these identifications
and determines $\mathcal O_K$ in general; the concrete-ring proofs here
do not require that later classification.
\item[$(\alpha),\ (\alpha,\beta)$.] Ideals generated by one or two elements:
the second consists of the sums $\alpha x+\beta y$ with
$x,y\in\mathcal O_K$. General ideal theory is developed in Volume~13.
\item[$N(I)$.] The positive norm of a non-zero ideal $I$, distinct from
the element norm used for Gaussian and Eisenstein integers. A field norm
of an element in a real quadratic field need not be positive.
\item[$|x|_p$.] The $p$-adic absolute value, defined by $|x|_p=p^{-v_p(x)}$
for non-zero rationals and $|0|_p=0$.
\item[$\Z_p,\ \Q_p$.] The $p$-adic integers and their field of fractions:
they are not the same as the finite residue ring $\Z/p\Z$.
\end{description}


\section*{From arithmetic to symmetry}
A number can reveal a larger structure when we change which operations
we study. The nonzero classes modulo $7$ multiply: $3,2,6,4,5,1$ are
the successive powers of $3$. This is a group with a generator, not
merely a list of remainders. The Gaussian integers contain $5=(2+i)(2-i)$:
they form a ring in which factorisation takes on new patterns. Finally,
the character $\chi_4$ satisfies $\chi_4(3)=-1$ and
$\chi_4(3\cdot3)=1=\chi_4(3)^2$; on the invertible classes its rule
preserves multiplication.

Try one final check before entering Volume~3. In $(\Z/4\Z)^\times$ the class of $3$ has square $1$, and $\chi_4(3)^2=(-1)^2=1=\chi_4(1)$. This is a concrete instance of a map preserving multiplication. The next book will study the homomorphism rule in much larger groups.

These three examples are the invitation to Volume~3,
\emph{Symmetry}. It develops groups, rings, fields and homomorphisms in
general, with the concrete arithmetic of this book now available as a
source of examples. Volume~6 will later ask how often primes occur, and
Volume~13 will repair unique factorisation with ideals in more general
number fields. Those are later questions. The next one is how to
recognise the structure common to all three examples.

\section*{Illustration credits}
The mathematical diagrams are original vector drawings for this volume. The following historical reproductions
are in the public domain and were obtained through Wikimedia Commons: Pierre de Fermat,
traditional portrait, artist unknown; Carl Friedrich Gauss, portrait by Christian Albrecht Jensen, 1840;
Pafnuty Chebyshev, photograph, late 1880s or early 1890s.
Their public-domain status is independent of the copyright in this book's text and new illustrations.
